\documentclass[11pt]{report}
\usepackage[utf8]{inputenc}
\usepackage[T1]{fontenc}
\usepackage{lmodern}
\usepackage[english]{babel}
\usepackage{graphicx}
\usepackage{hyperref}
\usepackage{amsmath}
\usepackage{xcolor}
\usepackage{listings}
\usepackage{longtable}
\usepackage{booktabs}
\usepackage{geometry}
\geometry{margin=0.85in}
\setlength{\parskip}{2pt}

% --- Code listing style ---------------------------------------------------
\definecolor{codebg}{rgb}{0.96,0.96,0.96}
\definecolor{codecomment}{rgb}{0.4,0.45,0.4}
\definecolor{codestring}{rgb}{0.5,0.1,0.1}
\lstdefinestyle{shell}{
    backgroundcolor=\color{codebg},
    basicstyle=\ttfamily\small,
    breaklines=true,
    frame=single,
    numbers=none,
    language=bash,
    commentstyle=\color{codecomment},
    keywordstyle=\bfseries,
    aboveskip=6pt,
    belowskip=6pt
}
\lstdefinestyle{code}{
    backgroundcolor=\color{codebg},
    basicstyle=\ttfamily\small,
    breaklines=true,
    frame=single,
    numbers=left,
    numberstyle=\tiny\color{gray},
    commentstyle=\color{codecomment},
    stringstyle=\color{codestring},
    keywordstyle=\bfseries,
    aboveskip=6pt,
    belowskip=6pt
}
\lstset{style=shell}

% --- "which machine" banner, used at the top of each chapter -------------
\newcommand{\runson}[1]{%
  \begin{center}
  \fcolorbox{black}{codebg}{\parbox{0.94\textwidth}{\textbf{Runs on:} #1}}
  \end{center}
}

% --- boxed warning, for facts that will cost the reader time if missed ----
\newsavebox{\warnbox}
\newenvironment{warningbox}%
  {\begin{lrbox}{\warnbox}\begin{minipage}{0.90\textwidth}}%
  {\end{minipage}\end{lrbox}%
   \begin{center}\fcolorbox{black}{codebg}{\usebox{\warnbox}}\end{center}}


\title{Carolus under ROS2 --- Lab PC and Raspberry Pi 5}
\author{}
\date{\today}

\begin{document}

\maketitle

\begin{abstract}
This document is deliberately \textbf{self-contained and independent of the
main Carolus/RoboMaster manual} (\texttt{overleaf/technical.tex}) and of the
rest of that project's repository. It, and the folder it lives in
(\texttt{raspberry5-carolus-ros2/}), can be copied, moved, or handed over on
their own and still work: the ROS2 wrapper source
(\texttt{carolus\_ros2/}) holds its own copy of the detection/solver core
rather than referencing another workspace by relative path, and every
install step below is written out in full rather than pointed at another
document. \textbf{The cost of that choice, stated plainly:} the C++ core
copied into \texttt{carolus\_ros2/src/} and \texttt{include/carolus\_node/}
is a snapshot, not a live link --- if the ROS1 node's own copy
(\texttt{carolus\_ws/src/libuvgs\_astrobee/}) is fixed or changed later, this
copy does not update itself, and someone has to notice and re-sync it by
hand. Accepted on purpose, because this document's job is to be portable,
not to be the single source of truth for that source code.

\textbf{What Carolus is, for a reader with no other context:} a physical
4-LED target (the \emph{beacon}) is detected in a camera frame, and a
Perspective-4-Point (P4P) solve on the four blob centres yields a full
6-DOF camera pose relative to it --- originally derived from NASA's SVGS
system for the Astrobee free-flyer, adapted here for a DJI RoboMaster S1.
The ROS1 node running this today is \texttt{carolus\_astrobee}.

This is the ROS2 side of Carolus, in full: what it took to build and run the
same detection/pose pipeline under ROS2, on two different machines. No
robot is involved anywhere in this document --- both targets run against a
plain USB webcam, or a synthetic image.
\end{abstract}

\tableofcontents

\chapter{Carolus under ROS2 --- Lab PC}
\label{ch:ros2-labpc}

\runson{the \textbf{lab PC} (or any x86\_64 machine with ROS2 installed). No
robot, no Pi, no RNDIS connection anywhere in this chapter.}

\section{Why this exists}

The Carolus/RoboMaster project's ROS1 node (\texttt{carolus\_astrobee}) mixed
its detection algorithm (blob preprocessing, contour extraction, blob
selection) with its ROS wiring inside one class. The detection methods'
\emph{signatures} were already framework-agnostic ---
\texttt{cv::Mat}/\texttt{std::vector}/\texttt{double}, no ROS types --- but
their \emph{bodies} called ROS logging macros directly, which is real
coupling a signature alone does not show.

That algorithm was extracted into a separate library, \texttt{carolus\_core}
(\texttt{beacon\_detector.cpp} plus \texttt{ceresP4P.cpp}/\texttt{pose\_est.cpp}/
\texttt{pose\_filter.cpp}), built with zero ROS headers and zero ROS linkage
--- confirmed by an empty \texttt{NEEDED} list on the compiled library, not
assumed from reading the source. The ROS1 node links against it instead of
containing it; this chapter (and the next) builds the exact same source a
second and third time, under different toolchains and middlewares, as the
actual test of that claim --- not that the code compiles, but that it
produces the same numbers.

\section{Installing ROS2 Humble}

Ubuntu 22.04 is Humble's own native target, so it installs from the official
packages:

\begin{lstlisting}[style=shell]
sudo apt install software-properties-common curl -y
sudo add-apt-repository universe -y
sudo curl -sSL https://raw.githubusercontent.com/ros/rosdistro/master/ros.key \
    -o /usr/share/keyrings/ros-archive-keyring.gpg
echo "deb [arch=$(dpkg --print-architecture) \
    signed-by=/usr/share/keyrings/ros-archive-keyring.gpg] \
    http://packages.ros.org/ros2/ubuntu $(. /etc/os-release && echo $UBUNTU_CODENAME) main" \
    | sudo tee /etc/apt/sources.list.d/ros2.list > /dev/null
sudo apt update
sudo apt install ros-humble-ros-base ros-humble-image-transport \
    ros-humble-cv-bridge python3-colcon-common-extensions \
    libeigen3-dev libceres-dev libopencv-dev build-essential cmake -y
\end{lstlisting}

\begin{warningbox}
\texttt{ros-humble-ros-base} alone does \textbf{not} include
\texttt{image\_transport} or \texttt{cv\_bridge} --- the same gap exists on
the ROS1 side of this project (\texttt{ros-noetic-ros-base} vs.
\texttt{ros-noetic-desktop-full}), and skipping either package here
reproduces the same broken-build failure on ROS2. The four packages after
\texttt{python3-colcon-common-extensions} are not ROS packages at all ---
they are \texttt{carolus\_core\_ros2}'s own C++ dependencies
(\texttt{find\_package(Eigen3/Ceres/OpenCV)} in its \texttt{CMakeLists.txt}),
easy to miss because this project's own lab PC already had them installed
from unrelated ROS1 work, which is exactly the trap: a dependency that is
satisfied by accident on the machine that wrote this document is invisible
until a genuinely blank machine tries to build it.
\end{warningbox}

\section{Installing ROS2 Jazzy, in a container}

Jazzy targets Ubuntu 24.04, not 22.04 --- installing its packages natively on
a 22.04 host risks library-version ABI mismatches between what Jazzy expects
and what the host actually has. A container sidesteps the question rather
than risking it:

\begin{lstlisting}[style=shell]
docker pull ros:jazzy-ros-base
docker run -d --name carolus_jazzy \
    -v /path/to/raspberry5-carolus-ros2:/carolus_ros2_ws \
    ros:jazzy-ros-base sleep infinity
docker exec carolus_jazzy apt-get update
docker exec carolus_jazzy apt-get install -y \
    ros-jazzy-image-transport ros-jazzy-cv-bridge \
    python3-colcon-common-extensions \
    libeigen3-dev libceres-dev libopencv-dev build-essential cmake
\end{lstlisting}

\section{The \texttt{carolus\_ros2} package}

\texttt{carolus\_ros2/} is a minimal \texttt{rclcpp} wrapper --- \textbf{a
proof that the core ports, not a full reimplementation of the ROS1 node.}
Unlike an earlier version of this package, it does \textbf{not} reference
another workspace by relative path: \texttt{src/} and
\texttt{include/carolus\_node/} hold their own copy of
\texttt{beacon\_detector}/\texttt{ceresP4P}/\texttt{pose\_est}/\texttt{pose\_filter},
so this whole folder builds on its own.

Deliberately out of scope, stated here so the gap is not mistaken for an
oversight later: the Astrobee-specific FOV undistortion path (plumb\_bob
only), \texttt{ff\_msgs} (dead code even on the ROS1 side --- the node
advertises a publisher on this message type and immediately overwrites it
with the real \texttt{/pose} publisher on the next line, so nothing has ever
consumed it), and the producer/consumer image queue (one synchronous
callback instead).

\begin{warningbox}
A ROS2 (\texttt{ament\_cmake}) package must \textbf{not} live inside a
catkin (ROS1) workspace's source space. \texttt{catkin\_make} refuses to
configure a workspace containing any non-catkin package --- \texttt{"This
workspace contains non-catkin packages in it, and catkin cannot build a
non-homogeneous workspace without isolation"} --- so placing one there
breaks the ROS1 build of \emph{every} package in that workspace, not just
its own. A \texttt{CATKIN\_IGNORE} marker does not solve it either:
\texttt{colcon} honours that file too, so the marker fixes the ROS1 build by
hiding the package from the ROS2 build. The two build systems need entirely
separate source trees, which is one reason this whole document and its
folder are kept independent.
\end{warningbox}

\begin{lstlisting}[style=shell]
# Host, Humble:
source /opt/ros/humble/setup.bash
mkdir -p ~/humble_ws/src && cd ~/humble_ws
ln -s /path/to/raspberry5-carolus-ros2/carolus_ros2 src/carolus_ros2
colcon build --packages-select carolus_ros2

# Inside the container (Jazzy), with this folder mounted at /carolus_ros2_ws:
source /opt/ros/jazzy/setup.bash
mkdir -p /colcon_ws/src
ln -s /carolus_ros2_ws/carolus_ros2 /colcon_ws/src/carolus_ros2
cd /colcon_ws && colcon build --packages-select carolus_ros2
\end{lstlisting}

\textbf{Verified 2026-08-18, Jazzy}: clean build, 11.3\,s, zero warnings
attributable to the new code. \texttt{ldd} on the resulting binary, run with
the ROS2 environment sourced, resolves every dependency.

\textbf{Verified 2026-08-18, Humble, same host}: clean build, 7.45\,s, same
zero-warning result. The test in the next section produced byte-identical
output to the Jazzy run above (same contour count, same solver iteration
count and cost, matching pose fields to 11 significant digits --- the
remaining digit is Ceres iteration-order floating-point noise, not a real
difference) --- the same source, unmodified, compiled by two different ROS2
toolchains, ran the same detection and solve correctly on both.

\begin{warningbox}
\texttt{cv\_bridge}'s header is mid-rename across ROS2 distros
(\texttt{cv\_bridge.h} deprecated, \texttt{cv\_bridge.hpp} its replacement)
--- and \emph{which one a given install actually has is not predictable
from the distro name alone}. This project's Jazzy container only has
\texttt{.hpp}; its Humble install (22.04) only has \texttt{.h} --- confirmed
directly with \texttt{dpkg -L ros-humble-cv-bridge | grep cv\_bridge}, not
assumed from a changelog. \texttt{carolus\_node\_ros2.cpp} includes
whichever header actually exists (\texttt{\#if
\_\_has\_include(<cv\_bridge/cv\_bridge.hpp>)}) rather than hardcoding one
distro's answer --- if porting this pattern elsewhere, check the installed
package directly on the target machine first.
\end{warningbox}

\section{Testing it, without a robot}

A synthetic beacon frame stands in for a live camera.

\begin{lstlisting}[style=shell]
sudo apt install -y python3-numpy python3-opencv
\end{lstlisting}

Four blobs, positioned by projecting the real beacon geometry through this
project's own camera intrinsics at a plausible working range, published on a
loop to the node's subscribed topic:

\begin{lstlisting}[style=code]
import numpy as np, cv2
fx, fy, cx, cy = 546.1957, 547.0838, 575.6041, 372.1876
Z = 0.7
points_3d = [(0.0825,0,0), (-0.0825,0,0), (0,0.072,0), (0,0,0.0555)]
# Pick the blob colour FROM HSV, not guessed in BGR -- a saturated
# hand-picked BGR value easily lands too dark to clear the preprocessing
# threshold once converted to grayscale (this is why: the saturation
# findAndCalcContours needs for its own hue mask pulls luminance down,
# working against the brightness the initial threshold wants).
hsv = np.uint8([[[115, 90, 255]]])   # hue in this project's blue window (90-140)
bgr = cv2.cvtColor(hsv, cv2.COLOR_HSV2BGR)[0][0].tolist()
img = np.zeros((720, 1280, 3), dtype=np.uint8)
for X, Y, _ in points_3d:
    u, v = int(fx*X/Z + cx), int(fy*Y/Z + cy)
    cv2.circle(img, (u, v), 14, bgr, -1)
cv2.imwrite('synth_beacon.png', img)
\end{lstlisting}

Publish it on \texttt{/camera/color/image\_raw} (any small
\texttt{rclpy} timer-based publisher using \texttt{cv\_bridge}), then:

\begin{lstlisting}[style=shell]
ros2 run carolus_ros2 carolus_node --ros-args -p image_threshold:=150
ros2 topic echo /pose --once
\end{lstlisting}

\begin{warningbox}
\texttt{image\_threshold:=150} is a test-only override for the synthetic
image above, not this project's real value (190, tuned for the physical
beacon). If a publisher process is already running, restarting the image
generator alone does \textbf{not} update what it sends --- a script that
loads the frame once at startup keeps serving the first version from memory
no matter how many times the file on disk changes. Restart the publisher,
not just the file.
\end{warningbox}

\subsection{QoS: the ROS2 failure that looks like success}

\begin{warningbox}
In ROS1 a subscriber connects to any publisher on the topic. In ROS2 it
connects only if the two QoS profiles are compatible, and an incompatible
pair delivers \textbf{nothing} --- no error, no warning, no log line. The
callback simply never fires and the node looks perfectly healthy. A
\textsc{reliable} subscriber specifically \emph{cannot} receive from a
\textsc{best effort} publisher --- the one QoS mismatch that silently drops
every frame rather than negotiating something that still works.
\end{warningbox}

\textbf{This specific driver does not trigger that trap, measured directly
rather than assumed:} \texttt{ros2 topic info /image\_raw --verbose} against
\texttt{ros-humble-usb-cam} 0.8.1 shows \texttt{Reliability: RELIABLE} ---
\emph{not} best effort. An earlier version of this section asserted "camera
drivers normally publish best effort" as a general claim; it was never
measured against this project's own driver, and the measurement contradicts
it. Both QoS profiles below were tested against the live C920 feed and both
received frames (\texttt{qos\_profile:=default}: 99 frames in 4\,s;
\texttt{qos\_profile:=sensor\_data}: 31 frames in 4\,s) --- verified
2026-08-20, not inferred from the general rule.

\textbf{The default is kept at \texttt{sensor\_data} anyway.} It is still
the strictly more compatible side --- a best-effort subscriber accepts
\emph{either} kind of publisher, so it can never be the thing that breaks
compatibility --- and the RoboMaster's own real camera bridge (not yet
tested through this ROS2 path) may turn out to publish best-effort, unlike
this one driver. \texttt{qos\_profile} stays overridable specifically so
that question can be tested directly against whatever publisher is in front
of it, rather than assumed either way.

\begin{center}
\begin{tabular}{@{}lll@{}}
\toprule
\textbf{\texttt{qos\_profile}} & \textbf{Reliability} & \textbf{Receives from} \\
\midrule
\texttt{sensor\_data} (default) & best effort & best-effort \emph{and} reliable publishers \\
\texttt{default}                & reliable    & reliable publishers only \\
\bottomrule
\end{tabular}
\end{center}

\begin{lstlisting}[style=shell]
# What a publisher actually offers, before assuming either QoS profile:
ros2 topic info /image_raw --verbose
\end{lstlisting}

\textbf{Verified 2026-08-18, Jazzy}: with a beacon-shaped frame actually
arriving, \texttt{findAndCalcContours} found exactly 4 contours,
\texttt{selectBlobs} returned a 4-blob group, target sorting succeeded, the
Ceres solver ran (31 iterations, did not converge on this synthetic scene ---
expected, the flat-depth projection above is not a fully self-consistent P4P
scene, and convergence was never the point), and a real
\texttt{geometry\_msgs/msg/PoseStamped} was published and read back,
correctly typed for ROS2 (\texttt{header.stamp.sec}/\texttt{.nanosec}, the
ROS2 field names --- ROS1 uses \texttt{.secs}/\texttt{.nsecs} on the same
field).

\section{Against a real camera: a USB webcam}
\label{sec:ros2-realcam}

\texttt{ros-humble-usb-cam} is a prebuilt binary --- no source build needed.
\texttt{rqt\_image\_view} is optional but useful here specifically: unlike
the synthetic test, what the camera actually sees is not known in advance.

\begin{lstlisting}[style=shell]
sudo apt install ros-humble-usb-cam ros-humble-rqt-image-view
\end{lstlisting}

\begin{lstlisting}[style=shell]
ros2 run usb_cam usb_cam_node_exe --ros-args \
    -p video_device:=/dev/video1 -p image_width:=1280 -p image_height:=720 \
    -p pixel_format:=yuyv2rgb
ros2 run rqt_image_view rqt_image_view /image_raw
\end{lstlisting}

\begin{warningbox}
\texttt{pixel\_format:=yuyv2rgb} is not optional on a device whose target
resolution offers no Motion-JPEG mode (\texttt{usb\_cam}'s own default,
\texttt{mjpeg2rgb}, silently falls back to raw \texttt{yuv422\_yuy2} in that
case --- check with the driver's own printed format list). A raw YUYV
buffer reaching \texttt{carolus\_node} crashes it: \texttt{findAndCalcContours}
calls into \texttt{cv::findContours}, which requires \texttt{CV\_8UC1}, and a
2-channel packed YUV image is neither that nor the 3-channel BGR the
callback otherwise assumes. \texttt{ros2 topic echo <topic> --field encoding
--once} shows what a driver is actually publishing when in doubt.
\end{warningbox}

\begin{lstlisting}[style=shell]
ros2 run carolus_ros2 carolus_node --ros-args \
    --params-file carolus_ros2/config/logitech_1080p.yaml
\end{lstlisting}

\texttt{logitech\_1080p.yaml} holds this camera's intrinsics
(fx=570.6975, fy=572.8081, cx=322.6828, cy=228.3519, distortion
[0.0223, -0.3088, -0.0048, -0.0018] --- a Logitech HD Pro Webcam C920, from
Hector directly, 2026-08-17, approximate, "should be calibrated again
later") and the same beacon geometry used elsewhere in this project.
\textbf{One field needs a real YAML list, not a string:}
\texttt{known\_points} must be \texttt{[0.0825, 0.0, 0.0, ...]}, a literal
list --- \texttt{declare\_parameter} has no string-to-vector parsing for a
\texttt{std::vector<double>} default, and a quoted string either fails to
load or silently mismatches the declared type.

\begin{warningbox}
\textbf{A second silent failure, this one undetectable from the log.}
\texttt{cv\_bridge::toCvShare(msg, msg->encoding)} returns whatever encoding
the source publishes --- against this driver, \texttt{rgb8}, not
\texttt{bgr8}. Converting RGB data with \texttt{cv::COLOR\_BGR2HSV} swaps red
and blue going into the hue calculation: no crash, no warning, every hue
silently rotated. A real blue-LED beacon can land clean outside
\texttt{lb\_hue}/\texttt{ub\_hue} as a direct result, and the node keeps
publishing \texttt{"0 contours found"} on schedule --- identical to the
correct output when no beacon is in frame. Fixed by requesting
\texttt{"bgr8"} explicitly from \texttt{cv\_bridge::toCvShare}, not the
source encoding: the conversion becomes \texttt{cv\_bridge}'s job for
whatever a given driver publishes, not an assumption baked into this
callback.
\end{warningbox}

\textbf{Verified 2026-08-20}: Logitech HD Pro Webcam C920, 1280x720, measured
10.000\,Hz average (std dev 0.00392\,s, 30-sample window --- this device
offers only 10\,FPS at this resolution, confirmed against its own format
list; 30\,FPS exists only at 640x480 and below). \texttt{carolus\_node}
processes every frame with no crash, and reports \texttt{"0 contours found"}
with no beacon in view --- the correct result, not a failure.

\textbf{With a lit beacon in frame, same session:} four blobs found,
\texttt{SortTargetsUsingTetrahedronGeometry} succeeded, the Ceres solve ran,
and a \texttt{geometry\_msgs/msg/PoseStamped} was published on \texttt{/pose}
with a correctly normalised quaternion. \textbf{One parameter change was
required to get there}, and it is not optional on this camera:

\begin{lstlisting}[style=shell]
ros2 run carolus_ros2 carolus_node --ros-args \
    --params-file carolus_ros2/config/logitech_1080p.yaml \
    -p image_threshold:=100
\end{lstlisting}

\begin{warningbox}
\textbf{A lit blue beacon can sit far below \texttt{image\_threshold} while
being perfectly correct in hue and saturation.} Measured on a real frame
grabbed off \texttt{/image\_raw}: the LEDs read hue 105--115 (well inside the
90--140 window) and saturation 200--255 (well above the 80 floor), but their
\emph{grayscale} value peaks at \textbf{131} --- under the default
\texttt{image\_threshold} of 190, so \texttt{preprocessImage} discards them
before the hue mask ever runs. Saturated blue converts to a low grayscale
value under the standard luminance weighting, where blue's coefficient is
smallest. The symptom is \texttt{"0 contours found"} on every frame with a
beacon plainly visible in \texttt{rqt\_image\_view}.

Do not guess a replacement value: grab one frame and measure. Moving the
beacon closer does \emph{not} fix it (tested --- still zero contours at
threshold 190), because the ceiling is luminance, not blob size.
\texttt{image\_threshold:=100} worked for this camera, this beacon and this
lighting; it is a per-setup value, not a new project default.
\end{warningbox}

\begin{warningbox}
\textbf{Detection working does not mean the pose is trustworthy.} Every solve
in this test reported \texttt{converged=0} (Ceres hit
\texttt{max\_num\_iterations=30} without converging), and the pose is written
\emph{unconditionally} --- a non-converged solve returns a finite, plausible
pose that is indistinguishable from a good one unless the diagnostic fields
are read. Two \texttt{/pose} reads taken at two visibly different beacon
positions returned translations a few centimetres apart
(\texttt{(4.232, 2.957, -7.476)} then \texttt{(4.201, 2.939, -7.444)}), which
is not what a solver tracking a moved target produces. Treat the pose values
from this setup as unvalidated: the detection path is confirmed, the pose
accuracy is not.
\end{warningbox}

\chapter{Carolus under ROS2 --- Raspberry Pi 5}
\label{ch:ros2-pi5}

\runson{a \textbf{Raspberry Pi 5}, no RoboMaster attached --- a plain USB
webcam (the same Logitech C920 as Chapter~\ref{ch:ros2-labpc}) is the only
sensor. \textbf{Everything below is planned, not yet verified on this board
--- see the status line at the start of each section.}}

\section{Why a Pi 5, and why aarch64 matters here}

Two things this chapter exists to test that the lab-PC chapter could not:
whether the extracted core (Chapter~\ref{ch:ros2-labpc}'s claim) is genuinely
architecture-independent, not just toolchain-independent --- everything
verified so far ran on x86\_64 --- and whether ROS2 Jazzy runs
\emph{natively}, without the Docker layer the lab PC needed because its host
OS is 22.04, not 24.04. The Pi 5 can run Ubuntu 24.04 directly, which is
Jazzy's own native target.

\section{Installing Ubuntu 24.04 and ROS2 Jazzy, natively}

\textbf{Status: planned, not yet run on this board.}

\begin{lstlisting}[style=shell]
# Ubuntu Server 24.04 LTS for Raspberry Pi, flashed via Raspberry Pi Imager
# or manually per Ubuntu's own arm64 Pi image instructions.

sudo apt install software-properties-common curl -y
sudo add-apt-repository universe -y
sudo curl -sSL https://raw.githubusercontent.com/ros/rosdistro/master/ros.key \
    -o /usr/share/keyrings/ros-archive-keyring.gpg
echo "deb [arch=$(dpkg --print-architecture) \
    signed-by=/usr/share/keyrings/ros-archive-keyring.gpg] \
    http://packages.ros.org/ros2/ubuntu $(. /etc/os-release && echo $UBUNTU_CODENAME) main" \
    | sudo tee /etc/apt/sources.list.d/ros2.list > /dev/null
sudo apt update
sudo apt install ros-jazzy-ros-base ros-jazzy-image-transport \
    ros-jazzy-cv-bridge ros-jazzy-usb-cam ros-jazzy-rqt-image-view \
    python3-colcon-common-extensions -y
\end{lstlisting}

\begin{warningbox}
\textbf{Not yet confirmed for this exact board.} The lab-PC chapter found
that \texttt{ros-*-ros-base} alone omits \texttt{image\_transport} and
\texttt{cv\_bridge} on Humble; assumed to hold for Jazzy on aarch64 too,
since it is the same buildfarm packaging convention, but not independently
checked here yet. Verify with \texttt{ros2 pkg list | grep -E
"image_transport|cv_bridge"} after install, before assuming the command
above was sufficient.
\end{warningbox}

\section{Building \texttt{carolus\_ros2} on aarch64}

\textbf{Status: planned, not yet run on this board.}

\begin{lstlisting}[style=shell]
source /opt/ros/jazzy/setup.bash
mkdir -p ~/pi5_ws/src && cd ~/pi5_ws
ln -s /path/to/raspberry5-carolus-ros2/carolus_ros2 src/carolus_ros2
colcon build --packages-select carolus_ros2
\end{lstlisting}

Expected to build unmodified --- \texttt{carolus\_core\_ros2}'s dependencies
(OpenCV, Eigen3, Ceres) are all available as \texttt{ros-jazzy-*}/plain
\texttt{apt} packages for arm64, and nothing in the source
(Chapter~\ref{ch:ros2-labpc}) uses an x86-specific intrinsic or a
compiler flag that would not carry over. \textbf{This is a prediction, not a
measurement} --- record the actual build time and any warning here once run,
the same way Chapter~\ref{ch:ros2-labpc} recorded 7.45\,s (Humble) and
11.3\,s (Jazzy) on the lab PC.

\section{Against the webcam}

\textbf{Status: planned, not yet run on this board.}

Same procedure as Section~\ref{sec:ros2-realcam}, same physical camera, same
config file (\texttt{carolus\_ros2/config/logitech\_1080p.yaml}) --- the
point of reusing it unmodified is that a difference in behaviour would then
be attributable to the board, not to a re-tuned parameter.

\begin{lstlisting}[style=shell]
ros2 run usb_cam usb_cam_node_exe --ros-args \
    -p video_device:=/dev/video0 -p image_width:=1280 -p image_height:=720 \
    -p pixel_format:=yuyv2rgb
ros2 run rqt_image_view rqt_image_view /image_raw
ros2 run carolus_ros2 carolus_node --ros-args \
    --params-file carolus_ros2/config/logitech_1080p.yaml \
    -p image_threshold:=100
\end{lstlisting}

\begin{warningbox}
\texttt{video\_device} is written as \texttt{/dev/video0} above, not the
lab PC's \texttt{/dev/video1} --- device enumeration is host-specific
(driven by what else is attached and in what order), never assume the same
device path carries over between machines. Confirm with
\texttt{v4l2-ctl --list-devices} (or, if unavailable,
\texttt{ls /dev/video*} plus \texttt{udevadm info}) before launching.
\end{warningbox}

\texttt{image\_threshold:=100} is carried over from the lab PC (see
Section~\ref{sec:ros2-realcam}'s first warningbox for why the default 190
finds nothing on this camera). It is a per-setup value --- if the Pi 5 run
finds nothing, re-measure on a frame from \emph{that} board's own feed rather
than assuming 100 transfers.

\textbf{Success criteria, not yet measured on this board.} Two stages, in
order --- do not conflate them:

\begin{enumerate}
  \item \textbf{Pipeline healthy, no beacon in frame}: \texttt{carolus\_node}
  runs with no crash, logs \texttt{"Time to find contours"} at the camera's
  actual rate, reports \texttt{"0 contours found"}. Record the measured rate.
  \item \textbf{Beacon in frame}: four blobs found, target sort succeeds,
  Ceres runs, \texttt{/pose} publishes. Record \texttt{final\_cost} and
  whether \texttt{converged} is 0 or 1 --- on the lab PC it was consistently
  0, and whether the Pi 5 behaves the same is one of the more interesting
  things this chapter can establish.
\end{enumerate}

Both stages passed on the lab PC (Section~\ref{sec:ros2-realcam}). Matching
stage 1 but failing stage 2 would point at the camera/lighting/threshold;
failing stage 1 would point at the board, the driver or the build.

\end{document}
